\chapter{Durchführung} \label{ch:Durchfuehrung}

Auch der Folgende Text ist am Praktikumsskript von Jens Wagner orientiert. Es gelten diesleben Hinweise wie in \cch{Einleitung}\cite{skriptPAP21,Duden}.
\section{Messverfahren} \label{sec:Messverfahren}
Die Durchführung gliedert sich in zwei Versuchteile, die sich auf verschiedenen Messmethoden beruhen.

\subsection{Versuch nach Clément-Desormes}
Zunächst wird der große Glasbehälter mit dem Luftbalg vorsichtig aufgepumpt, bis das Flüssigkeitsmanometer eine Höhendifferenz $h_1$ anzeigt. Es ist essentiell zu warten, bis das Gas die durch die Kompression entstandene Wärme an die Umgebung abgegeben hat (isothermer Zustand). Dies ist erreicht, wenn der Flüssigkeitsspiegel über ca. ein bis zwei Minuten konstant bleibt. 
Anschließend wird der Hahn kurzzeitig geöffnet und sofort wieder geschlossen, sobald der Druckausgleich mit der Atmosphäre (hörbares Zischen stoppt) erfolgt ist. Das Gas im Inneren ist nun kälter als die Umgebung. Man wartet erneut die isochore Erwärmung ab, bis die Manometerflüssigkeit ihren maximalen Endstand $h_3$ erreicht hat. Dieser Vorgang wird zur Fehlerminimierung mehrfach wiederholt.

\subsection{Versuch nach Rüchardt}
Für die dynamische Messung wird das Glasrohr vertikal auf den Behälter gesetzt. Zuerst wird die Messung mit Luft durchgeführt. Über ein Nadelventil wird ein feiner Gasstrom eingeleitet, der den Schwingkörper (den präzisen Metallzylinder) zum Schweben bringt. Durch eine kleine Entlüftungsöffnung im Rohr wird der Körper zu Schwingungen angeregt. 
Sobald eine stabile Schwingung ohne Wandberührung vorliegt, wird mit der Stoppuhr die Zeit für 50 Perioden gemessen. Dieser Messvorgang wird für Luft dreimal wiederholt. Im Anschluss wird das System gründlich mit Argon gespült, um sicherzustellen, dass sich kein Stickstoff oder Sauerstoff mehr im Behälter befindet, und die Messung der Schwingungsdauer wird analog für das Edelgas durchgeführt. Während der gesamten Durchführung ist auf die Sauberkeit des Glasrohrs zu achten, um Reibungseffekte zu minimieren.
